\ClearShipoutPicture 
\AddToShipoutPicture{%
 	\AtPageLowerLeft{\includegraphics[width=\paperwidth,height=\paperheight,page=4]{../notext_lb.pdf}}
}
	
	% Header
	\placetextbox{0}{0.935}{\paperwidth}{\centering \addfontfeature{LetterSpace=0}
\textrm{\fontsize{26}{1} \ding{66}}
{\fontsize{30}{1}\scshape Naomi Bishop}
\textrm{\fontsize{26}{1} \ding{66}}}
	\placetextbox{0}{0.9}{\paperwidth}{\fontsize{14}{1}\centering\itshape\addfontfeature{LetterSpace=0} Ehemalige Grubenkämpferin und Leibwächterin von Lady Blackbird
}
	
	% Traits
	\placetextbox{0.089}{0.86}{0.4\paperwidth}{\fontsize{13}{0}\fontcelestiaantiqua\scshape EIGENSCHAFTEN}
	\placetextbox{0.089}{0.83}{0.4\paperwidth}{\fontsize{12}{0}\fontcelestiaantiqua
Grubenkämpfer \vspace{38pt} \\
Leibwache \vspace{26pt} \\
Ehemalige Sklavin \vspace{26pt} \\
Scharfsinnig}
	\placetextbox{0.12}{0.813}{0.37\paperwidth}{\fontsize{10}{0}\fontcelestiaantiqua\itshape
Kampferfahren, Brutal, Wandelnde Waffe, Schnell, Abgehärtet, [Stark], [Knochenbrecher], [Furchteinflößendes Aussehen] \vspace{17pt} \\
Achtsamkeit, Bedrohungen, Verteidigen, Entwaffnen, Festnehmen, Tragen, Verzögern, [Sicherheit], [Erste Hilfe] \vspace{17pt} \\
Schleichen, Verstecken, Rennen, Robust, Aushalten, Betteln, Adel, [Hass], [Eiserner Wille] \vspace{17pt} \\
Verständnisvoll, Aufmerksam, Startbereit, Lügner, Fallen, [Gefahren], [Motive Erkennen] }

	% Keys and Secrets
	\placetextbox{0.515}{0.865}{0.4\paperwidth}{\fontsize{12}{0}\fontcelestiaantiqua
Schlüssel des Wächters \vspace{50pt} \\
Schlüssel der Rache \vspace{66pt} \\
Schlüssel des Kriegers \vspace{52pt} \\
Geheimnis der Zerstörung \vspace{30pt} \\
Geheimnis der Leibwache}
	\placetextbox{0.545}{0.848}{0.37\paperwidth}{\fontsize{10}{0}\fontcelestiaantiqua
\textit{Du bist das loyale Schild von Lady Blackbird. Benutze diesen Schlüssel, wenn sie deine Entscheidungen beeinflusst oder wenn du sie vor Gefahren beschützt.} Austausch: Beende dein Verhältnis zu Lady Blackbird. \vspace{15pt} \\
\textit{Das Imperium hielt dich als Sklaven und hat dich gezwungen, für ihr Vergnügen zu töten. Du wirst deine Rache bekommen und ihre Welten in Flammen sehen. Benutze diesen Schlüssel, wenn du dem Imperium Schaden zufügst; insbesondere, wenn du einen Imperialen tötest.}
Austausch: Vergib dem Imperium, was es dir angetan hat. \vspace{15pt} \\
\textit{Du suchst Auseinandersetzungen und den Blutrausch des Kampfes; je härter, desto besser. Benutze diesen Schlüssel, wenn du einen Kampf gegen einen würdigen oder überlegenen Gegner führst.}
Austausch: Lasse die Chance auf einen guten Kampf verstreichen. \vspace{14pt} \\
\textit{Du kannst Dinge mit deinen bloßen Händen auseinandernehmen wie andere Menschen mit einem Vorschlaghammer. Furchteinflößend.} \vspace{11pt} \\
\textit{Einmal pro Sitzung kannst du einen Würfelwurf wiederholen, der genutzt wurde um jemanden zu beschützen.} }
	
	
	%
	% Rules Summary 1
	\placetextbox{0.089}{0.435}{0.4\paperwidth}{\fontsize{8}{0}\fontcelestiaantiqua 
{\fontsize{14}{0}\fontcelestiaantiqua\scshape Rolling the Dice} \vspace{3pt} \\
When you try to overcome an obstacle, you roll dice. Start with one die.
Add a die if you have a \textbf{trait} that can help you. If that trait has any \textbf{tags}
that apply, add another die for each tag. Finally, add any number of dice
from your personal \textbf{pool} of dice (your pool starts with 7 dice). \vspace{5pt} \\
Roll all the dice you’ve gathered. Each die that shows \textbf{4 or higher} is a hit.
You need hits equal to the difficulty \textbf{level} (usually 3) to pass the obstacle. \vspace{5pt} \\
	\textsc{levels: 2 easy—3 difficult—4 challenging—5 extreme} \vspace{5pt} \\
	\textbf{If you pass,} discard all the dice you rolled (including any pool dice you
used). Don’t worry, you can get your pool dice back. \vspace{5pt} \\
	\textbf{If you don't pass,}  you don’t yet achieve your goal. But, you get to keep
the pool dice you rolled and \textbf{add another die to your pool}. The GM will
escalate the situation in some way and you might be able to try again. \vspace{6pt} \\
{\fontsize{14}{0}\fontcelestiaantiqua\scshape Conditions} \vspace{3pt} \\
	When events warrant or especially when you fail a roll, the GM may
impose a \textbf{condition} on your character: \textbf{Injured, Dead, Tired, Angry,
Lost, Hunted,} or \textbf{Trapped}. When you take a condition, mark its box
and say how it comes about. [\textsc{Note}: The “dead” condition just means
“presumed dead” unless you say otherwise.] \vspace{6pt} \\
{\fontsize{14}{0}\fontcelestiaantiqua\scshape Helping} \vspace{3pt} \\
	If your character is in a position to help another character, you can give them
a die from your pool. Say what your character does to help. If the roll
fails, you get your pool die back. If it succeeds, your die is lost.
	}
	
	% Rules Summary 2
	\placetextbox{0.515}{0.435}{0.4\paperwidth}{\fontsize{8}{0}\fontcelestiaantiqua 
{\fontsize{14}{0}\fontcelestiaantiqua\scshape Keys} \vspace{3pt} \\
When you hit a Key, you can do one of two things:
\begin{itemize}
	\setlength\itemsep{0em}
	\item Take an \textbf{experience point} (XP)
	\item Add a die to your pool (up to a max of 10)
\end{itemize}
	If you go into danger because of your key, you get 2 XP or 2 pool dice
(or 1 XP and 1 pool die). When you have accumulated 5 XP, you earn an
\textbf{advance}. You can spend an advance on one of the following:
\begin{itemize}
	\setlength\itemsep{0em}
	\item Add a new \textbf{Trait} (based on something you learned during play or
on some past experience that has come to light)
	\item Add a \textbf{tag} to an existing trait
	\item Add a new \textbf{Key} (you can never have the same key twice)
	\item Learn a \textbf{Secret} (if you have the means to)
\end{itemize}
	You can hold on to advances if you want, and spend them at any time,
even in the middle of a battle! \vspace{5pt} \\
	Each key also has a \textbf{buyoff}. If the buyoff condition occurs, you have the
option of removing the Key and earning two advances. \vspace{6pt} \\
{\fontsize{14}{0}\fontcelestiaantiqua\scshape Refresh} \vspace{3pt} \\
	You can refresh your pool back to 7 dice by having a \textbf{refreshment scene}
with another character. You may also remove a condition or regain the
use of a Secret, depending on the details of the scene. A refreshment
scene is a good time to ask questions (in character) so the other player
can show off aspects of his or her PC—“Why did you choose this
life?”—“What do you think of the Lady?”—“Why did you take this
job?” etc. Refreshment scenes can be flashbacks, too.
	}	
	
	\placetextbox{0.089}{0.555}{0.4\paperwidth}{\fontsize{9}{0}\fontcelestiaantiqua  \addfontfeature{LetterSpace=0}
Schlagwörter in [Klammern] sind Fähigkeiten, die du noch nicht erlangt hast. Du kannst sie bei einem Fortschritt freischalten. Lies genaueres in der Regelzusammenfassung.}
	
\placetextbox{0.155}{0.506}{0.1\paperwidth}{\fontsize{12}{0}\fontcelestiaantiqua\addfontfeature{LetterSpace=-5}\scshape verletzt}
\placetextbox{0.281}{0.508}{0.1\paperwidth}{\fontsize{12}{0}\fontcelestiaantiqua\addfontfeature{LetterSpace=-5}\scshape tot}
\placetextbox{0.376}{0.508}{0.1\paperwidth}{\fontsize{12}{0}\fontcelestiaantiqua\addfontfeature{LetterSpace=-5}\scshape wütend}
\placetextbox{0.478}{0.508}{0.1\paperwidth}{\fontsize{12}{0}\fontcelestiaantiqua\addfontfeature{LetterSpace=-5}\scshape verirrt}
\placetextbox{0.587}{0.508}{0.1\paperwidth}{\fontsize{12}{0}\fontcelestiaantiqua\addfontfeature{LetterSpace=-5}\scshape müde}
\placetextbox{0.681}{0.508}{0.1\paperwidth}{\fontsize{12}{0}\fontcelestiaantiqua\addfontfeature{LetterSpace=-5}\scshape verfolgt}
\placetextbox{0.802}{0.508}{0.1\paperwidth}{\fontsize{12}{0}\fontcelestiaantiqua\addfontfeature{LetterSpace=-5}\scshape gefangen}